\section{Methodology}
\label{sec:method}

\subsection{Task and Dataset}
\label{sec:dataset}

We study short fiction generation (400--600 words) using prompts from the \writingprompts dataset~\citep{fan2018writingprompts}, a large corpus of creative writing prompts sourced from Reddit.
We randomly sample 15 diverse prompts (seed $= 42$) from a pool of 5{,}000, covering science fiction, horror, fantasy, literary fiction, and humor.
Prompt lengths range from 30 to 300 characters.
While 15 prompts is a modest sample, it is sufficient for detecting the large effects we observe, and we report effect sizes alongside significance tests to contextualize the results.

\subsection{Experimental Conditions}
\label{sec:conditions}

We compare six conditions, summarized in \tabref{tab:conditions}:

\begin{table}[t]
\centering
\caption{Experimental conditions. API calls are counted per prompt.}
\label{tab:conditions}
\begin{tabular}{@{}llc@{}}
\toprule
\textbf{Condition} & \textbf{Description} & \textbf{API Calls} \\
\midrule
\baseline & Direct generation, no critique & 1 \\
\selfrefine $\times$1 & Generate $\rightarrow$ Critique $\rightarrow$ Revise & 3 \\
\selfrefine $\times$2 & Two \critiquerevise iterations & 5 \\
\selfrefine $\times$3 & Three \critiquerevise iterations & 7 \\
\bestofn ($N\!=\!4$) & Generate 4 stories, self-rank, select best & 5 \\
\bestofn + Refine & \bestofn, then one \critiquerevise pass & 7 \\
\bottomrule
\end{tabular}
\end{table}

\para{Self-Refine conditions.}
Following \citet{madaan2023selfrefine}, each iteration consists of three steps: (1)~generate or take the current story, (2)~produce a structured critique identifying specific strengths and weaknesses along six dimensions, and (3)~revise the story addressing the critique.
We test one to three iterations to measure whether improvement is monotonic.

\para{Best-of-N conditions.}
To disentangle the effect of critique from mere generation diversity, we generate four independent stories per prompt and have the model rank them.
We also test combining selection with one refinement pass.

\subsection{Model Configuration}
\label{sec:config}

We use \gptfour for all roles (generation, critique, and evaluation) with role-specific temperature settings to separate the capacities (\tabref{tab:config}).
All generation uses seed $= 42$; judging uses seed $= 1042$.

\begin{table}[t]
\centering
\caption{Model configuration by role. All roles use \gptfour.}
\label{tab:config}
\begin{tabular}{@{}lccc@{}}
\toprule
\textbf{Parameter} & \textbf{Generation} & \textbf{Critique} & \textbf{Judging} \\
\midrule
Temperature & 0.8 & 0.4 & 0.3 \\
Max tokens & 800 & 500 & 600 \\
Seed & 42 & 42 & 1042 \\
\bottomrule
\end{tabular}
\end{table}

\subsection{Evaluation Protocol}
\label{sec:evaluation}

\para{Blind pairwise comparison.}
For each of the 15 prompts, we compare every non-baseline condition against the baseline in a blind pairwise evaluation.
The judge receives two stories (labeled ``Story A'' and ``Story B'') with no information about which condition produced them.
Presentation order is randomized per comparison to control for position bias.

\para{Multi-dimensional scoring.}
The judge scores each story on six HANNA-inspired dimensions~\citep{chhun2022hanna}, each on a 1--5 scale:
\begin{itemize}[leftmargin=*,itemsep=0pt,topsep=0pt]
    \item \textbf{Coherence}: Plot logic, internal consistency, narrative flow
    \item \textbf{Creativity}: Originality of ideas, premise, narrative choices
    \item \textbf{Engagement}: How compelling and readable the story is
    \item \textbf{Character}: Depth, believability, distinctiveness of characters
    \item \textbf{Prose Quality}: Writing craft, imagery, word choice, voice
    \item \textbf{Emotional Impact}: Ability to evoke genuine feelings
\end{itemize}
The judge also declares an overall winner or tie.

\para{Reward hacking indicators.}
To detect superficial quality inflation, we measure word count, vocabulary richness (type-token ratio), and bigram diversity (Distinct-2) across conditions.

\subsection{Statistical Analysis}
\label{sec:stats}

We use Wilcoxon signed-rank tests for pairwise comparisons of average dimension scores, with Bonferroni correction for five comparisons ($\alpha = 0.01$).
We report Cohen's $d$ as an effect size measure.
The total experiment comprises 90 generated stories (15 prompts $\times$ 6 conditions) and 75 pairwise evaluations (15 prompts $\times$ 5 non-baseline conditions).
